% Generated from validated Article Draft Result. Do not hand-edit; revise the structured draft instead.
\section{動画・画像生成が研究デモから競争軸へ}
\noindent\textbf{2024年上期、生成AIの競争はテキストと画像理解だけでなく、動画・高品質画像・音声を「生成する」側へ急速に広がった。Sora、GoogleのVeo/Imagen 3、Stable Diffusion 3、Klingを中心に、能力そのものと公開形態を分けて追う。}\autocite{src-0f33134a9129,src-0fe343211040,src-7b3484ead314,src-9c32b95d27a7}

Soraの発表は、text-to-videoを研究デモの周辺領域からfrontier model competitionの中心へ押し出した。同じ上期にはGoogleがVeoやImagen 3を含むgenerative mediaをI/Oで提示し、KuaishouはKlingを発表、RunwayもGen-3 Alphaを打ち出した。動画生成は短いclipの noveltyではなく、各社がmodel familyとproduct strategyを結び付ける競争軸になった。\autocite{src-0f33134a9129,src-0fe343211040,src-9c32b95d27a7,src-b881d1321d94}


画像側ではStable Diffusion 3が2月のpreviewから6月のMedium releaseへ進み、研究発表、API、weight公開の時点差を持つlifecycleを示した。media modelでは「発表された」ことと「一般ユーザーが使える」「APIで利用できる」「weightを取得できる」ことが大きく異なり、この区別なしに競争を比較すると時系列を誤る。\autocite{src-7b3484ead314}


同時期のGPT-4oやChameleonはmultimodal understanding/generationの境界を別方向から動かした。入力として画像や音声を理解する能力と、高品質な動画・画像を新規生成する能力は別の技術軸である。2024-H1には両者が同時進行したため、「multimodal」という一語でまとめず、interaction modelとmedia generatorを分離して読む必要がある。\autocite{src-0f33134a9129,src-0fe343211040,src-2d3a533e91db,src-e1259b3c0804}


\begin{claimboundary}
各社が公開したsample、benchmark、品質比較は、それぞれの選定条件と評価方法に依存する。preview映像の存在は一般提供や再現可能性を意味せず、API提供はopen-weight releaseとも異なる。本稿では能力claimとavailability stateを別々に記述する。\autocite{src-0f33134a9129,src-0fe343211040,src-36cdd427b766,src-7b3484ead314,src-9c32b95d27a7}
\end{claimboundary}

この半年の意味は、生成メディアが「面白い出力例」から、制作workflow、広告、映像、design toolへ接続し得るplatform capabilityとして認識され始めたことにある。ただしH1時点では公開範囲や利用条件に大きな差が残り、技術的潜在力と実用可能性の距離そのものが競争の一部だった。\autocite{src-0f33134a9129,src-0fe343211040,src-7b3484ead314,src-9c32b95d27a7,src-b881d1321d94}
